\documentclass[tikz,border=6pt]{standalone}
% Shared visual language for the thesis architecture figures.
% Compile each standalone source with XeLaTeX from this directory.
\usepackage[UTF8,scheme=plain,fontset=none]{ctex}
\usepackage{amsmath,amssymb}
\usepackage{fontspec}
\setmainfont{TeX Gyre Termes}
\setsansfont{TeX Gyre Heros}
\setmonofont{DejaVu Sans Mono}
\setCJKmainfont{Noto Serif SC}
\setCJKsansfont[BoldFont=SimHei]{Noto Sans SC}
\usetikzlibrary{arrows.meta,positioning,calc,fit,backgrounds,shapes.geometric,matrix,decorations.pathreplacing}

\definecolor{Ink}{HTML}{203040}
\definecolor{Muted}{HTML}{66788A}
\definecolor{LineSoft}{HTML}{AAB7C4}
\definecolor{DataBlue}{HTML}{2563A6}
\definecolor{DataFill}{HTML}{EAF3FB}
\definecolor{ClockPurple}{HTML}{7450A5}
\definecolor{ClockFill}{HTML}{F1ECF8}
\definecolor{ControlTeal}{HTML}{16827B}
\definecolor{ControlFill}{HTML}{E9F6F4}
\definecolor{DspPurple}{HTML}{6A4C93}
\definecolor{DspFill}{HTML}{F0EAF7}
\definecolor{RamGreen}{HTML}{367C59}
\definecolor{RamFill}{HTML}{ECF6F0}
\definecolor{FabricOrange}{HTML}{C56A1A}
\definecolor{FabricFill}{HTML}{FFF1E5}
\definecolor{EvidenceGreen}{HTML}{2F7D4B}
\definecolor{EvidenceFill}{HTML}{EAF6EE}
\definecolor{Amber}{HTML}{B07A16}
\definecolor{AmberFill}{HTML}{FFF6DD}
\definecolor{PendingGray}{HTML}{7C8792}
\definecolor{PendingFill}{HTML}{F1F3F5}
\definecolor{Danger}{HTML}{A04444}
\definecolor{DangerFill}{HTML}{FFF0F0}

\tikzset{
  >={Stealth[length=2.2mm,width=1.45mm]},
  every node/.style={font=\sffamily\footnotesize,text=Ink},
  data/.style={-{Stealth[length=2.2mm,width=1.45mm]},draw=DataBlue,line width=0.92pt},
  data bi/.style={{Stealth[length=1.9mm,width=1.3mm]}-{Stealth[length=1.9mm,width=1.3mm]},draw=DataBlue,line width=0.82pt},
  control/.style={-{Stealth[length=2.0mm,width=1.35mm]},draw=ControlTeal,dashed,line width=0.78pt},
  clock/.style={-{Stealth[length=2.0mm,width=1.35mm]},draw=ClockPurple,densely dotted,line width=0.9pt},
  address/.style={-{Stealth[length=2.0mm,width=1.35mm]},draw=FabricOrange,dash dot,line width=0.76pt},
  feedback/.style={-{Stealth[length=1.8mm,width=1.25mm]},draw=Muted,line width=0.66pt},
  block/.style={draw=Ink,fill=white,rounded corners=1.5pt,line width=0.68pt,minimum height=10mm,minimum width=23mm,align=center,inner sep=3pt},
  data block/.style={block,draw=DataBlue,fill=DataFill},
  control block/.style={block,draw=ControlTeal,fill=ControlFill},
  clock block/.style={block,draw=ClockPurple,fill=ClockFill},
  dsp/.style={block,draw=DspPurple,fill=DspFill,line width=0.9pt},
  ram/.style={block,draw=RamGreen,fill=RamFill,line width=0.9pt},
  fabric/.style={block,draw=FabricOrange,fill=FabricFill,line width=0.9pt},
  result/.style={block,draw=EvidenceGreen,fill=EvidenceFill,line width=0.9pt},
  pending/.style={block,draw=PendingGray,fill=PendingFill,line width=0.8pt},
  warning/.style={block,draw=Amber,fill=AmberFill,line width=0.9pt},
  boundary/.style={draw=Muted!75,rounded corners=2.2mm,dashed,line width=0.55pt,inner sep=3mm},
  title/.style={font=\sffamily\bfseries\large,text=Ink,align=center},
  subtitle/.style={font=\sffamily\scriptsize,text=Muted,align=center},
  group title/.style={font=\sffamily\bfseries\footnotesize,text=Ink,anchor=west},
  signal/.style={font=\sffamily\scriptsize,text=Ink,fill=white,inner sep=1.2pt,align=center},
  note/.style={font=\sffamily\scriptsize,text=Muted,align=center},
  tiny note/.style={font=\sffamily\tiny,text=Muted,align=center},
  badge/.style={draw=EvidenceGreen,fill=EvidenceFill,rounded corners=1.8mm,line width=0.65pt,font=\sffamily\bfseries\scriptsize,text=EvidenceGreen,inner xsep=4pt,inner ysep=2pt,align=center},
  pending badge/.style={draw=PendingGray,fill=PendingFill,rounded corners=1.8mm,line width=0.65pt,font=\sffamily\bfseries\scriptsize,text=PendingGray,inner xsep=4pt,inner ysep=2pt,align=center},
  innovation/.style={circle,draw=Amber,fill=AmberFill,line width=0.8pt,minimum size=6.5mm,font=\sffamily\bfseries\scriptsize,text=Amber,inner sep=0pt},
  sum/.style={circle,draw=Ink,fill=white,line width=0.75pt,minimum size=6mm,inner sep=0pt},
  mux/.style={trapezium,trapezium left angle=72,trapezium right angle=108,draw=Ink,fill=white,line width=0.68pt,minimum height=11mm,minimum width=12mm,align=center,inner sep=2pt}
}

\newcommand{\bitrate}[2]{{\scriptsize #1 bit，#2}}
\newcommand{\passmark}{\textcolor{EvidenceGreen}{\bfseries 满足判据}}
\newcommand{\pendingmark}{\textcolor{PendingGray}{\bfseries 待验证}}


\begin{document}
\begin{tikzpicture}[x=1cm,y=1cm]

\node[title] at (12.5,11.75) {基于时分MAC与双口BRAM的跨级资源复用微体系结构};
\node[subtitle] at (12.5,11.30)
  {当前2-DSP核心：Stage1独占一条串行MAC通道，Stage2/Stage3共享另一条MAC通道与统一存储体系};

% Unified coefficient memory at the top.
\node[ram,minimum width=81mm,minimum height=20mm] (coeffram) at (13.05,9.60)
  {RAMB18E1\#3：统一FIR系数双口存储器\\[-1pt]
   {\scriptsize Port A：地址128--179，Stage1展开52项，signed 16 bit}\\[-1pt]
   {\scriptsize Port B：地址0--63为Stage2/Stage3-flat；96--127为Stage3-P3补偿，signed 18 bit}};

% Stage1 lane.
\node[data block,minimum width=20mm] (x1) at (0.90,6.25) {$x[n]$\\[-1pt]{\scriptsize 24 bit}};
\node[ram,minimum width=31mm,minimum height=19mm] (h1) at (3.35,6.25)
  {RAMB18E1\#1\\[-1pt]{\scriptsize Stage1历史}\\[-1pt]{\scriptsize 64$\times$24 SDP；有效52字}};
\node[fabric,minimum width=24mm,minimum height=19mm] (scan1) at (6.15,6.25)
  {历史/系数地址扫描\\[-1pt]{\scriptsize 0--51逐项读取}\\[-1pt]{\scriptsize 中心延时索引25}};
\node[dsp,minimum width=31mm,minimum height=22mm] (dsp1) at (9.25,6.25)
  {DSP48E1\#0\\[-1pt]{\scriptsize 24$\times$16 + P串行MAC}\\[-1pt]{\scriptsize 52项MAC；含BRAM返回与尾拍流水}};
\node[data block,minimum width=18mm] (y2) at (11.80,6.25) {$y_2$\\[-1pt]{\scriptsize 24 bit}};
\draw[data] (x1) -- (h1);
\draw[data] (h1) -- (scan1);
\draw[data] (scan1) -- node[signal,above] {历史样值} (dsp1);
\draw[data] (dsp1) -- (y2);
\draw[address] (scan1.north) |- node[signal,pos=0.72,above] {Port A地址} (coeffram.west);
\draw[data] (coeffram.south west) -| node[signal,pos=0.70,left] {系数A} (dsp1.north);

% Shared Stage2/Stage3 lane.
\node[control block,minimum width=29mm,minimum height=18mm] (pending23) at (13.55,3.95)
  {$S_2/S_3$ pending\\[-1pt]{\scriptsize phase与补偿模式锁存}\\[-1pt]{\scriptsize 优先级：$S_2\rightarrow S_3$}};
\node[ram,minimum width=35mm,minimum height=20mm] (h23) at (16.95,3.95)
  {RAMB18E1\#2\\[-1pt]{\scriptsize Stage2/3统一历史}\\[-1pt]{\scriptsize 5 bit地址=bank+index}\\[-1pt]{\scriptsize 32个逻辑字，signed 20 bit}};
\node[fabric,minimum width=30mm,minimum height=20mm] (sched) at (20.30,3.95)
  {任务调度与输入MUX\\[-1pt]{\scriptsize job\_stage / phase / index}\\[-1pt]{\scriptsize $S_2$: 9/8 MAC；$S_3$: 6/5 MAC}};
\node[dsp,minimum width=32mm,minimum height=22mm] (dsp23) at (23.55,3.95)
  {DSP48E1\#1\\[-1pt]{\scriptsize 20$\times$18 + P共享MAC}\\[-1pt]{\scriptsize Q15舍入、PATTERN溢出检测}\\[-1pt]{\scriptsize 结果按job\_stage解复用}};
\node[data block,minimum width=19mm] (s2in) at (13.55,6.25) {$S_2$桥接输入\\[-1pt]{\scriptsize 20 bit}};
\node[data block,minimum width=19mm] (s3in) at (16.10,6.25) {$S_3$桥接输入\\[-1pt]{\scriptsize 20 bit}};
\node[data block,minimum width=19mm] (y4) at (22.50,7.05) {$y_4$\\[-1pt]{\scriptsize 20 bit}};
\node[data block,minimum width=19mm] (y8) at (24.40,7.05) {$y_8$\\[-1pt]{\scriptsize 21 bit}};
\draw[control] (s2in.south) -- (pending23.north west);
\draw[control] (s3in.south) -- (pending23.north east);
\draw[address] (pending23) -- (h23);
\draw[data] (h23) -- (sched);
\draw[data] (sched) -- (dsp23);
\draw[data] (dsp23.north west) |- (y4.south);
\draw[data] (dsp23.north east) |- (y8.south);
\draw[address] (sched.north) |- node[signal,pos=0.68,right] {Port B地址} (coeffram.east);
\draw[data] (coeffram.south east) -- ++(0,-3.25)
  -- node[signal,above] {系数B} (23.55,5.35) -- (dsp23.north);

% Board-only memory and boundary cards.
\node[pending,minimum width=35mm,minimum height=17mm] (rom) at (2.25,2.20)
  {板级RAMB18E1\#4\\[-1pt]{\scriptsize 256$\times$32双速测试音ROM}\\[-1pt]{\scriptsize 不属于滤波器核心}};
\node[result,minimum width=47mm,minimum height=18mm] (total) at (7.50,2.20)
  {核心物理专用资源\\[-1pt]{\bfseries 2 DSP48E1 + 3 RAMB18E1}\\[-1pt]{\scriptsize 整板另含1个测试音RAMB18E1}};

% Innovation markers.
\node[innovation] (i1) at (2.00,8.35) {I1};
\node[note,anchor=west,align=left] at (2.40,8.35)
  {26个对称系数展开为52项，\\避免预加器并顺序复用单DSP};
\node[innovation] (i2) at (16.50,8.35) {I2};
\node[note,anchor=west,align=left] at (16.90,8.35)
  {Stage2/3跨级任务调度，\\共享MAC及尾拍定点处理};
\node[innovation] (i3) at (12.60,1.05) {I3};
\node[note,anchor=west] at (13.00,1.05) {两级历史逻辑分区、物理统一BRAM};
\node[innovation] (i4) at (17.70,1.05) {I4};
\node[note,anchor=west] at (18.10,1.05) {三级系数与P3补偿库统一寻址};
\node[innovation] (i5) at (22.65,1.05) {I5};
\node[note,anchor=west,align=left] at (23.05,1.05) {PREG使能屏蔽\\未初始化历史抽头};

% Background grouping.
\begin{scope}[on background layer]
  \node[boundary,fill=DataFill!32,fit=(x1)(h1)(scan1)(dsp1)(y2)(i1)] (lane1) {};
  \node[boundary,fill=ControlFill!28,fit=(s2in)(s3in)(pending23)(h23)(sched)(dsp23)(y4)(y8)(i2)] (lane23) {};
\end{scope}
\node[group title] at ($(lane1.north west)+(0.18,-0.25)$) {通道A：Stage1串行MAC};
\node[group title] at ($(lane23.north west)+(0.18,-0.25)$) {通道B：Stage2/Stage3时分共享MAC};

\node[note,anchor=west] at (0.25,0.15)
  {实线为样值/系数数据，青色虚线为任务控制，橙色点划线为地址与调度；RAM共享不改变各逻辑级的独立状态语义。};

\end{tikzpicture}
\end{document}


